\documentclass[a4paper,12pt]{report}
\usepackage[T1]{fontenc} 
\usepackage[brazil]{babel}    %%% Idioma default brazil
\usepackage[utf8]{inputenc}
\usepackage{ae}
\usepackage{graphicx,url,geometry}
\usepackage{amsmath,amssymb}
\usepackage{float}   
\usepackage{color}
\usepackage{booktabs}
\usepackage{multirow}
\usepackage{steinmetz}
\usepackage{rotating}
\usepackage{fancyhdr}



\geometry{a4paper, left=3cm, right=2cm, top=2cm, bottom=2cm}



\begin{document}

\chapter*{Malha de aterramento}

Usando como caso de estudo um solo estratificado cuja resistividade decresce linearmente com a profundidade até que se mantém constante em 4 metros, fizemos os cálculos para uma malha de terra supondo que a área do terreno em que o prédio está localizado é de $20\times20$ metros. Consideramos uma corrente de curto de 10kA.

\begin{table}[htbp]
  \centering

    \begin{tabular}{cc}
    \toprule
    \textbf{$A_{(m)}$} & \textbf{$\rho_{(\Omega)}$}\\
    \midrule
    1     & 400 \\ 
    2     & 200 \\
    4     & 100 \\
    \bottomrule
    \end{tabular}%
  \caption{Resistividade usada nos cálculos}
\end{table}

\begin{itemize}
\item ${\rho}_{1}$ - 400 $\Omega.m$     
\item ${\rho}_{2}$ - 100 $\Omega.m$   
\item $\frac{\rho_2}{\rho_1}= 0,25 \longrightarrow K_1=0,7981 $ 
\item $\rho_m = K_1\times\rho_1 = 319,24$ $\Omega.m$ 
\item $R_{(raio, m)} = \sqrt{S/\pi} = \sqrt{120/\pi} = 11,3m$
\item $S = 20\times20 = 400m^2$
\item $\rho_{(H=1m)} = 400\Omega ; \rho_{(H=4m)}=100\Omega $ Interpolando:
\item $\rho_{(H=2,96m)} = \rho_m = 319,24$ $\Omega.m$ $\longrightarrow D_p = 2,96m $
\item $R/D_p = 11,3/2,96 = 6,251m = K_2 \longrightarrow$ Interpolando: $K_{3(\rho_2/\rho_1, K_2)} \tilde{=} 0,8$
\item $\rho_a = K_3\times\rho_1 = 240$ $\Omega.m$ 
\item $S_{(condutor)} = 0,002533\times 10kA = 25,3mm^2$ adotamos $S_c=25mm^2$ para cabo simples com solda exotérmica e duração de falha $T_f = 0,5s$
\item Primeira tentativa: $D_l = 0,2\times C_m $; $D_c = 0,2\times L_m \longrightarrow N_{cp}= 5 + 1 = 6$; $N_{cj} = 5 + 1 = 6$
\item Comprimento dos condutores: $L_{cm}= 1,05\times[(C_m\times N_{cj})+(L_m\times N_{cp})] = 252m$
\item $R_{mc} = \frac{\rho_a}{4\times R} + \frac{\rho_a}{L_{cm}} = 6,26\Omega$, que cumpre o requerimento $R_{mc}<10\Omega$
\item Resistência de um eletrodo de terra $R_{le} = \frac{\rho_a}{2\times\pi\times L_h}\times\ln(\frac{400\times L_h}{2,54\times D_h}) $ estipulando $D_h = 3/4"$ e $L_h=2,4m $(norma) $\longrightarrow R_{le} = 99,03\Omega$
\item $K_h = \frac{1+A\times B}{N_h} $ $A=0,1629 ;B=14,06 $ portanto $K_h=0,063215$ e $R_{ne}=6,3\Omega$ 
\item Resistência Mútua $R_{mu} = \frac{\rho_m}{\pi\times L_{cm}}\times[\ln(\frac{2\times L_{cm}}{L_{th}}+K_1\times\frac{L_{cm}}{\sqrt{S}} - K_2 + 1)] = 3,31\Omega$
\item Resistência total da malha: $R_{tm} = \frac{R_{mc}\times R_{ne} - R_{mu}^2}{R_{mc}+R_{ne}-2\times R_{mu}} = 0,26\Omega < 10\Omega$

\end{itemize}

\end{document}